% Part: methods
% Chapter: induction
% Section: relations

\documentclass[../../../include/open-logic-section]{subfiles}

\begin{document}

\olfileid{mth}{ind}{rel}

\olsection{关系与函数}

当我们归纳地定义了一类对象（比如自然数或者良构项）之后，我们也能够归纳地定义关于这些对象的\emph{关系}。例如，考虑这样的想法：如果一个良构项~$t_1$ 作为另一个良构项~$t_2$ 的一部分出现，则称 $t_1$ 是 $t_2$ 的\textit{子项}。为此引入一个符号：$t_1 \sqsubseteq t_2$。当然，每个良构项都是自身的子项：$t \sqsubseteq t$。我们可以如下归纳地定义这个关系：

\begin{defn}
  良构项~$t_1$ 作为~$t_2$ 的子项的关系 $t_1 \sqsubseteq t_2$，通过对~$t_2$ 归纳定义如下：
  \begin{enumerate}
  \item 若 $t_2$ 是一个字母，则 $t_1 \sqsubseteq t_2$ 当且仅当 $t_1 = t_2$。
  \item 若 $t_2$ 是 $[s_1 \circ s_2]$，则 $t_1 \sqsubseteq t_2$ 当且仅当 $t_1 = t_2$，$t_1 \sqsubseteq s_1$，或 $t_1 \sqsubseteq s_2$。
  \end{enumerate}
\end{defn}

例如，这个定义将告诉我们 $\mathrm{a}
\sqsubseteq [\mathrm{b} \circ \mathrm{a}]$。根据 (2)，$\mathrm{a} \sqsubseteq [\mathrm{b} \circ \mathrm{a}]$ 当且仅当 $\mathrm{a}
= [\mathrm{b} \circ \mathrm{a}]$，或 $\mathrm{a} \sqsubseteq b$，或 $\mathrm{a} \sqsubseteq \mathrm{a}$。前两个为假：$\mathrm{a}$ 显然不等于 $[\mathrm{b} \circ
  \mathrm{a}]$，而根据 (1)，$\mathrm{a} \sqsubseteq \mathrm{b}$ 当且仅当 $\mathrm{a} = \mathrm{b}$，这也是假的。然而，同样根据 (1)，$\mathrm{a} \sqsubseteq \mathrm{a}$ 当且仅当 $\mathrm{a} = \mathrm{a}$，这是真的。

值得注意的是，这个定义的成功依赖于一个我们尚未证明的事实：每个良构项~$t$ 要么是一个字母，要么存在\emph{唯一确定}的良构项 $s_1$ 和 $s_2$，使得 $t = [s_1 \circ s_2]$。这里的“唯一确定”意味着，如果 $t = [s_1 \circ s_2]$，它就不会\emph{同时也}等于某个满足 $s_1 \neq r_1$ 或 $s_2 \neq r_2$ 的 $[r_1 \circ r_2]$。如果出现这种情况，那么第 (2) 条可能会自相矛盾：将 $t_2$ 读作 $[s_1 \circ s_2]$ 时我们可能得到 $t_1 \sqsubseteq t_2$，但如果将 $t_2$ 读作 $[r_1 \circ r_2]$ 我们可能得到并非 $t_1 \sqsubseteq t_2$。在我们证明这不可能发生之前，先看一个它\emph{可能}发生的例子。

\begin{defn}
  按如下方式归纳地定义\emph{无括号项}：
  \begin{enumerate}
  \item 每个字母都是一个无括号项。
    \item 如果 $s_1$ 和 $s_2$ 是无括号项，那么 $s_1 \circ s_2$ 也是一个无括号项。
    \item 除此之外都不是无括号项。
  \end{enumerate}
\end{defn}

无括号项的例子有：$\mathrm{a}$，$\mathrm{b} \circ \mathrm{d}$，$\mathrm{b} \circ \mathrm{a} \circ \mathrm{b}$。现在，如果我们像上面那样为无括号项定义“子项”，第二条将变为：

\begin{center}
  如果 $t_2 = s_1 \circ s_2$，则 $t_1 \sqsubseteq t_2$ 当且仅当 $t_1 = t_2$, $t_1 \sqsubseteq s_1$，或 $t_1 \sqsubseteq s_2$。
\end{center}

现在，$\mathrm{b} \circ \mathrm{a} \circ \mathrm{b}$ 具有 $s_1 \circ s_2$ 的形式，其中
\begin{align*}
  s_1 & = \mathrm{b} \text{ 且} & s_2 & = \mathrm{a} \circ \mathrm{b}.
\intertext{它也具有 $r_1 \circ r_2$ 的形式，其中}
  r_1 & = \mathrm{b} \circ \mathrm{a} \text{ 且} & r_2 &= \mathrm{b}.
\end{align*}
那么 $\mathrm{a} \circ \mathrm{b}$ 是 $\mathrm{b} \circ
\mathrm{a} \circ \mathrm{b}$ 的子项吗？按第一种解读，答案是“是”；按第二种解读，答案则是“否”。

良构项由其他良构项构建的方式是唯一的，这一性质称为\emph{唯一可读性}。由于在这类归纳定义的对象上归纳地定义关系很重要，我们必须证明它成立。

\begin{prop}
  假设 $t$ 是一个良构项。那么 $t$ 要么是一个字母，要么存在唯一确定的良构项 $s_1$, $s_2$，使得~$t =
  [s_1 \circ s_2]$。
\end{prop}

\begin{proof}
  如果 $t$ 是一个字母，则条件已满足。因此假设 $t$ 不是一个字母。从归纳定义我们可以看出，此时 $t$ 必定具有 $[s_1 \circ s_2]$ 的形式，其中 $s_1$ 和 $s_2$ 是某些良构项。剩下的就是要证明它们是唯一确定的，即如果 $t = [r_1 \circ r_2]$，那么 $s_1 = r_1$ 且 $s_2 = r_2$。

  于是假设 $t = [s_1 \circ s_2]$ 并且对于良构项 $s_1$, $s_2$, $r_1$, $r_2$ 也有 $t = [r_1 \circ r_2]$。我们需要证明 $s_1 = r_1$ 且 $s_2 = r_2$。
  首先，$s_1$ 和 $r_1$ 必须相等，否则其中一个将是另一个的真初始段。但根据 \olref[sti]{prop:initial}，如果 $s_1$ 和 $r_1$ 都是良构项，这是不可能的。而如果 $s_1 = r_1$，那么显然也有 $s_2 = r_2$。
\end{proof}

我们也可以归纳地定义函数：例如，我们可以定义函数~$f$，它将任意良构项映射到其中嵌套 $[\dots]$ 的最大深度，定义如下：

\begin{defn}
  \ollabel{defn:depth} 良构项的\emph{深度} $f(t)$ 归纳定义如下：
  \[
  f(t) = \begin{cases}
    0 & \text{ 若 $t$ 是一个字母}\\
    \max(f(s_1), f(s_2)) + 1 & \text{ 若 $t = [s_1 \circ s_2]$。}
  \end{cases}
  \]
\end{defn}

例如
\begin{align*}
  f([\mathrm{a} \circ \mathrm{b}]) & = 
    \max(f(\mathrm{a}),f(\mathrm{b})) + 1 = \\ 
  &= \max(0, 0) + 1 = 1, \text{ 及}\\
  f([[\mathrm{a} \circ \mathrm{b}] \circ \mathrm{c}]) & = 
    \max(f([\mathrm{a} \circ \mathrm{b}]), f(\mathrm{c})) + 1 = \\ 
  & = \max(1,0) + 1 = 2.
\end{align*}

这里，我们当然假设 $s_1$ 和 $s_2$ 是良构项，并且利用了以下事实：每个良构项要么是一个字母，要么具有 $[s_1 \circ s_2]$ 的形式。同样重要的是，它只能以一种方式具有这种形式。为了明白为什么，再次考虑我们之前定义的无括号项。对应的“定义”会是：
\[
  g(t) = 
  \begin{cases}
    0 & \text{ 若 $t$ 是字母}\\
   \max(g(s_1), g(s_2)) + 1 & \text{ 若 $t = s_1 \circ s_2$。}
  \end{cases}
\]
现在考虑无括号项 $\mathrm{a} \circ \mathrm{b} \circ
\mathrm{c} \circ \mathrm{d}$。它可以以不止一种方式解读，例如，解读为 $s_1 \circ s_2$，其中
\begin{align*}
  s_1 & = \mathrm{a} \text{ 与} & 
  s_2 & = \mathrm{b} \circ \mathrm{c} \circ \mathrm{d},
\intertext{或解读为 $r_1 \circ r_2$，其中}
  r_1 & = \mathrm{a} \circ b \text{ 与} & 
  r_2 &= \mathrm{c} \circ \mathrm{d}.
\end{align*}
按第一种解读计算 $g$ 会得到
\begin{align*}
  g(s_1 \circ s_2) & =
  \max(g(\mathrm{a}), g(\mathrm{b} \circ \mathrm{c} \circ \mathrm{d})) + 1 =\\ & =
  \max(0,2) + 1 = 3
  \intertext{而按照另一种解读我们得到}
  g(r_1 \circ r_2) & =
  \max(g(\mathrm{a} \circ \mathrm{b}), g(\mathrm{c} \circ \mathrm{d})) + 1=\\ &
  = \max(1,1) + 1 = 2
\end{align*}
但是一个函数必须总是给出唯一的值；因此我们对~$g$ 的“定义”根本没有定义出一个函数。

\begin{prob}
  给出函数 $l$ 的归纳定义，其中 $l(t)$ 是良构项~$t$ 中的符号个数。
\end{prob}

\begin{prob}
  对良构项 $t$ 进行结构归纳，证明 $f(t) < l(t)$（其中 $l(t)$ 是~$t$ 中的符号个数，而 $f(t)$ 是 $t$ 的深度，定义见 \olref[mth][ind][rel]{defn:depth}）。
\end{prob}

\end{document}